\chapter{支持向量机：用间隔约束分类}
\label{ch:svm}

支持向量机（Support Vector Machine, SVM）在分类得分中同时考虑间隔与违例代价。
先理解线性情形，再理解核函数如何改变相似性计算。

\begin{keyidea}[得分不直接等于概率]
线性得分$s=\vct w^\top\vct x+b$的符号给出二分类规则。
间隔关注样本离决策边界的程度；若需要概率，应另按适当流程估计或校准。
\end{keyidea}

\section{硬间隔与支持向量}
\label{sec:svm-margin}
训练中两类均出现且数据线性可分时，硬间隔形式最小化$\frac12\lVert\vct w\rVert^2$，
约束$y_i(\vct w^\top\vct x_i+b)\geq1$，其中$y_i\in\{-1,1\}$。
在这个规范化下，两条间隔边界为$s=\pm1$，总几何宽度为$2/\lVert\vct w\rVert$。

\begin{figure}[htbp]
\centering\begin{tikzpicture}
\begin{axis}[width=121mm,height=57mm,xmin=-2.6,xmax=2.6,ymin=-1.1,ymax=1.8,
 xlabel={$x_1$},ylabel={$x_2$},xtick={-2,-1,0,1,2},ytick={0,1},axis lines=left]
\path[fill=Pale] (axis cs:-1,-1.1)--(axis cs:1,-1.1)--(axis cs:1,1.8)--(axis cs:-1,1.8)--cycle;
\addplot[Teal,very thick] coordinates {(0,-1.1)(0,1.8)};
\addplot[Quiet,dashed] coordinates {(-1,-1.1)(-1,1.8)};
\addplot[Quiet,dashed] coordinates {(1,-1.1)(1,1.8)};
\addplot[only marks,mark=o,Teal,mark size=3pt] coordinates {(-2,1)(-1,0)};
\addplot[only marks,mark=triangle*,Amber,mark size=3pt] coordinates {(1,0)(2,1)};
\addplot[only marks,mark=o,draw=Ink,mark size=6pt] coordinates {(-1,0)(1,0)};
\draw[<->,Amber,thick] (axis cs:-1,-.55)--(axis cs:1,-.55);
\node[lab,below] at (axis cs:0,-.6) {总间隔宽度2};
\node[lab] at (axis cs:0,1.55) {$s=x_1=0$};
\end{axis}
\end{tikzpicture}

\caption{人工例子$\vct w=(1,0)$、$b=0$，圆形为负类，三角为正类；圈出的点位于间隔边界。}
\par\smallskip
\readfigure{决策边界$x_1=0$，间隔边界为$x_1=\pm1$，总宽度为2。
在对偶表示中，系数$\alpha_i>0$的训练点称支持向量；它们在硬间隔解中位于间隔边界。}
\end{figure}

\quick{硬间隔要求可分且对异常记录敏感；不可分时不能机械强制所有约束。}

\section{软间隔：正确分类也可能付出代价}
\label{sec:svm-soft}
软间隔允许松弛变量$\xi_i\geq0$：
\[
 \min_{\vct w,b,\xi}\frac12\lVert\vct w\rVert^2+C\sum_i\xi_i,
 \qquad y_is_i\geq1-\xi_i,\quad C>0.
\]
最小可行松弛为$\xi_i=\max(0,1-y_is_i)$，即Hinge损失。
本式使用损失求和，软件或其他章节若改为平均，参数缩放需相应比较。

\begin{center}
\begin{tabularx}{\textwidth}{P{27mm}P{24mm}Y}
\toprule
正类样本得分$s$ & 最小$\xi$ & 解释\\
\midrule
1.2 & 0 & 正确且在正侧间隔边界之外。\\
0.4 & 0.6 & 正确，但进入间隔内部。\\
0 & 1 & 在决策边界上，类别取决于平票约定。\\
$-0.4$ & 1.4 & 落到错误一侧。\\
\bottomrule
\end{tabularx}
\end{center}

$\xi=0$不只表示“恰在间隔边界”；$\xi>0$也不都表示误分类。
$C$越大，单位违例的惩罚越强，模型通常更努力适应训练约束；这不保证泛化更好。
支持向量也可能在间隔内部或被误分，不能仅按“离边界最近”一句话定义全部软间隔情形。

\begin{pitfall}[尺度会改变几何含义]
线性间隔和常见核函数都受输入尺度影响。标准化应在训练内拟合；
类别不平衡还可能需要类权重与合适评估，而不是只调整$C$。
\end{pitfall}

\section{核函数：在相似性上建立非线性边界}
\label{sec:svm-kernel}
核函数满足$K(x,z)=\langle\phi(x),\phi(z)\rangle$，
允许在不显式构造全部$\phi$坐标的情况下计算相应内积。
核SVM得分可写为$s(x)=\sum_i\alpha_i y_iK(x_i,x)+b$。
常见RBF核为
\[
 K(x,z)=\exp[-\gamma\lVert x-z\rVert_2^2],\qquad\gamma>0.
\]

\begin{figure}[htbp]
\centering\begin{tikzpicture}
\begin{axis}[width=122mm,height=50mm,xmin=0,xmax=3,ymin=0,ymax=1.05,
 xlabel={距离 $d=\lVert x-z\rVert_2$},ylabel={$K(x,z)$},
 legend style={font=\small,at={(.98,.98)},anchor=north east}]
\addplot[Teal,thick,domain=0:3,samples=90]{exp(-.25*x*x)};\addlegendentry{$\gamma=0.25$}
\addplot[Amber,thick,domain=0:3,samples=90]{exp(-x*x)};\addlegendentry{$\gamma=1$}
\end{axis}
\end{tikzpicture}

\caption{RBF相似度随距离降低；较大$\gamma$使影响更局部。}
\par\smallskip
\readfigure{相同点的核值为1。$x=(0,0)$、$z=(1,1)$、$\gamma=0.5$时，核值为$e^{-1}\approx0.368$。
指数的负号不能省略。核方法也不保证任意数据与标签都能理想泛化。}
\end{figure}

$C$与$\gamma$通常需要联合验证。核矩阵与求解成本可能随样本数迅速增加，
较大任务可比较线性方法、核近似等方案。核相似性、得分和概率各有含义，不能混用。

\subsection{自测与回看}
最小可行松弛$\xi=0$是否只表示样本恰在间隔边界？\textbf{不是，间隔之外也可以为0。}
正类得分$s=0.4$时，是否已经分错？\textbf{没有；但仍有0.6的间隔违例。}

\begingroup\small
\noindent\textbf{参考：}ISLR第二版第9章；scikit-learn Support Vector Machines文档。
图和数值例子为本书独立构造。
\par\endgroup
